%% ma: L10-B14-0037
%% lop: 10
%% chuong: 5
%% bai: 14
%% dang: 10.14.3
%% loai: TLN
%% muc: VD
%% dap_an: "4"
%% nguon: "chuyen-de-cac-so-dac-trung-cua-mau-so-lieu-khong-ghep-nhom-toan-10-chuong-trinh-moi.pdf (GV. Phan Nhật Linh) (P04), Bài 03 Dạng 1 Phần III Câu 4"
%% kien_thuc: "Lớp 10 Bài 13 (tứ phân vị), Bài 14 (khoảng tứ phân vị, giá trị bất thường)"
%% trang_thai: chinh-thuc
%% ghi_chu: "Tự duyệt theo yêu cầu giáo viên 22/10/2026."
\begin{ex}
Có bao nhiêu giá trị bất thường trong mẫu số liệu sau?
\begin{center}
\begin{tabular}{|c|c|c|c|c|c|c|c|c|c|}\hline
$5$ & $6$ & $19$ & $21$ & $22$ & $23$ & $24$ & $25$ & $26$ & $27$\\\hline
$28$ & $29$ & $30$ & $31$ & $32$ & $33$ & $34$ & $48$ & $49$ & \\\hline
\end{tabular}
\end{center}
\shortans{4}
\loigiai{$n=19$: $Q_2=x_{10}=27$; $Q_1=x_5=22$; $Q_3=x_{15}=32$; $\Delta_Q=10$.

$Q_1-1{,}5\Delta_Q=7$; $Q_3+1{,}5\Delta_Q=47$.

Các giá trị bất thường: $5;6;48;49$. Có $4$ giá trị.}
\end{ex}
